% Copyright (c) 2026 Eugene Joseph M. Ragasa.
% SPDX-License-Identifier: Apache-2.0
% Author: Eugene Joseph M. Ragasa
% Formal proof completed: 2026-10-27
\documentclass[11pt]{article}
\usepackage[T1]{fontenc}\usepackage{amsmath,amssymb}\usepackage[margin=1in]{geometry}\usepackage[hidelinks]{hyperref}
\title{Proof Dossier: Scalar-Conjugation Fourier Invariance}\author{Eugene Joseph M. Ragasa}\date{Formal proof completed: 2026-10-27}
\begin{document}\maketitle
\section{Identity and status}
ID \texttt{BG-FOURIER-THM-050}; status \texttt{proved}. Lean declaration:
\par\noindent{\tiny\path{ProjectKoios.BandGauge.Fourier.}}\\
{\tiny\path{finiteFourierCoefficient_scalarPhaseConjugation}}.
\section{Statement and proof}
Every finite Fourier coefficient of a scalar operator field is unchanged by arbitrary
pointwise unit-modulus scalar conjugation. The proof substitutes the pointwise identity
from \texttt{BG-FOURIER-THM-049}.
\section{Boundary}
This theorem must not be used as a higher-rank or vector-state gauge statement. Kernel
axioms are \texttt{propext}, \texttt{Classical.choice}, and \texttt{Quot.sound}; no
admissions. Checked with Lean \texttt{4.35.0} and Mathlib revision
\texttt{56a005855ea17e0b0958f73ec801ae140b0a1679}. AI-assisted and kernel-checked;
no independent expert review recorded.
\end{document}
